% BEGIN_REPORT::RESEARCH::RSR-0004::2026-10-07T18:26:59-07:00
\documentclass[11pt]{article}
\usepackage[T1]{fontenc}
\usepackage[margin=1in]{geometry}
\usepackage{amsmath,amssymb,amsthm,hyperref}
\emergencystretch=3em
\newtheorem{theorem}{Theorem}
\newtheorem{lemma}{Lemma}
\newcommand{\SAT}{\operatorname{SAT}}
\title{RSR-0004: q-Horn as a Strict Zero-Child Portfolio Member\\
and the Missing Endpoint-Compatible Space Bound}
\author{RESEARCH, ARC/WRAPPER research team}
\date{7 October 2026}
\begin{document}
\maketitle

\begin{abstract}
The primary q-Horn literature gives linear-time recognition and SAT decision
and shows that q-Horn properly contains Horn, Krom/2-CNF, and renamable Horn.
This transfers exactly to a zero-child partial WRAPPER member and eliminates
duplicate treatment of those tractable islands.  A bounded audit finds a
five-variable q-Horn formula outside Horn, dual-Horn, 2-CNF, renamable Horn,
the all-true/all-false phases, and size-two strong Horn backdoors.  The
transfer nevertheless does not close WRAPPER: the safe explicit-graph ledger
uses linear writable space, giving exponent sum two, while the Williams
endpoint is $2\cos(\pi/7)<2$.  No cited theorem supplies the missing joint
time--space implementation.  General favorable coverage, ARC closure, and
P versus NP remain open.
\end{abstract}

\section{Claim card and selected arrow}
\begin{description}
\item[] \textbf{Claim IDs:} QHORN-VALUE1 and QHORN-SPACE-GATE1.
\item[] \textbf{Status before run:} q-Horn was present in historical controls
and Drive portfolio records but absent as a charged member of the current
WRAPPER spine.  Its endpoint-compatible realization was \textsc{Open}.
\item[] \textbf{Desired conclusion:} determine exactly what the q-Horn
literature transfers to the current forward-only portfolio and what resource
obligation remains.
\item[] \textbf{Exact quantifiers:} there exists one fixed uniform partial
machine $W_q$ such that for every explicit CNF $F$, $W_q(F)$ either returns
\textsc{fail}, or returns the exact bit $\SAT(F)$; it succeeds on every
q-Horn $F$.  No assertion is made on failure.
\item[] \textbf{Input/model:} ordinary explicitly indexed CNF of encoded
length $M$, read-only input, random-access writable work memory.  Time and
space belong to the same realization.
\item[] \textbf{SAME\_REALIZATION:} the actual q-Horn recognizer/decider run
that returns the Boolean VALUE.  A separately postulated low-space solver is
not substituted into its time ledger.
\item[] \textbf{Weakest useful theorem:} an exact domain enlargement plus a
proof that the published linear-time result alone does not certify
$C+\sigma<2\cos(\pi/7)$.
\item[] \textbf{What would falsify the transfer:} a q-Horn formula on which
the declared certifying-function test or exact base-class decider is
incorrect.  What would close the remaining gate is a model-matched q-Horn
implementation with a proved exponent pair below the endpoint.
\end{description}

\section{Corpus reviewed}
The repository preflight passed.  I surveyed the current repository tree and
read root \texttt{main.tex}, \texttt{main2.tex}, \texttt{main3.tex}, the
verified current \texttt{main4.tex} and \texttt{main5.tex},
\texttt{AGENTS.md}, \texttt{CODEX.md}, the complete pnp-research skill and
its research-protocol and experiment-mode references, the proof-spines
protocol, \texttt{spines.bib}, all seven task specifications, and the
materially relevant code and raw results.  The immediately preceding PROOF
report
\path{research/agents/proof/2026-10-07_181151_PROOF_WRP-0004.tex}
was fetched from GitHub at blob
\texttt{9225f4b47d189520acb3b88118564d7aa847ef1d}; it was substantive and
contained its final verification sentinel, so no write recovery was needed.

Relevant current reports read at artifact level included RSR-0002,
RED-0002, RSR-0003, LEAD ARC-0004, and WRP-0004.  GitHub search also exposed
the q-Horn localization sections in \texttt{main2.tex}/\texttt{main3.tex}
and the exact-2P2N records.  The historical Google Drive document
\emph{q-Horn Portfolio Domination and Heterogeneous Backdoor-DNF Corollary}
(id \texttt{1s8tfk9BVdtBv\_sIkianMM141YI1CEA35-zTipuaQPGo}) was read in
full rather than through its status label.  It already proves, conditional on
its separate exact-2P2N radius premise, that treating Horn, dual-Horn, Krom,
and renamable Horn as independent base-class coordinates duplicates q-Horn.

\section{Duplication and contradiction audit}
\textbf{Duplicate components adopted.}  The Drive record already states the
class containments and warns that formula-class structure does not transfer
to arbitrary-machine hardness.  RSR-0002 and RED-0002 already charge
backdoor acquisition and show that decision-only VALUE members may bypass
large explicit residuals.  WRP-0004 already provides all-true/all-false phase
VALUE returns.  None of these is claimed as new.

\textbf{New reconciliation.}  The current WRAPPER implementation listed
Horn and 2-CNF VALUE procedures separately but did not state the single
q-Horn member or its work-space ledger.  QHORN-VALUE1 consolidates and
strictly extends that finite domain.  QHORN-SPACE-GATE1 prevents the
linear-time literature label from being silently promoted to a favorable
joint time--space point.

\textbf{Contradiction resolved.}  There is no conflict between the Drive
q-Horn radius obstruction on a special exact-2P2N family and the present
positive q-Horn solver.  The former is a representation/backdoor statement;
the latter directly decides formulas already in the tractable base class.
Neither is an arbitrary SAT-machine lower bound.

\section{Primary literature and exact transfer}
Boros, Crama, and Hammer introduced the generalized Horn class and its
polynomial inference properties \cite{boros1990qh_inference}.  Boros,
Hammer, and Sun give linear-time recognition of q-Horn formulas
\cite{boros1994qhorn}.  Gaspers et al. restate the exact certificate:
$F$ is q-Horn when there is
\[
 \beta:\operatorname{lit}(F)\to\{0,\tfrac12,1\},\qquad
 \beta(x)=1-\beta(\neg x),
\]
with
\[
 \sum_{\ell\in C}\beta(\ell)\le1
 \quad\text{for every clause }C,
\]
and record that both recognition and SAT decision are linear time and that
q-Horn properly contains Horn, Krom, and renamable Horn
\cite{gaspers2013qhorn}.

\begin{theorem}[QHORN-VALUE1]
There is a fixed uniform exact zero-child partial wrapper which, on every
explicit CNF $F$ of length $M$, returns \textsc{fail} when recognition fails
and otherwise returns the exact Boolean value $\SAT(F)$.  It uses no call to
the hypothetical general SAT decider and has $O(M)$ time in the cited q-Horn
model.
\end{theorem}
\begin{proof}
Run the linear-time q-Horn recognizer.  If it rejects, make no semantic claim
and return \textsc{fail}.  If it accepts, run the linear-time exact q-Horn
SAT algorithm and return its bit.  Thus all returned bits are exact and no
recursive child exists.  The same procedure covers Horn, dual-Horn after
global complementation, 2-CNF, and renamable Horn by class containment.
\end{proof}

This is a \textsc{Known Theorem} transfer at the base-class interface and a
\textsc{Proved/Exact} portfolio composition using the already proved
zero-child contract.  It does not detect a q-Horn backdoor in an arbitrary
formula, construct a fractional-size SAT child, or cover every CNF.

\section{Space gate and model matching}
The recognition proof uses a linear-size quadratic cover and implication
graph.  Explicit construction plus graph state gives the conservative
simultaneous ledger
\[
 T(M)=M^{1+o(1)},\qquad S(M)=M^{1+o(1)}.
\]

\begin{lemma}[QHORN-SPACE-GATE1]
The cited linear-time q-Horn theorem alone does not establish an
endpoint-compatible WRAPPER member.  Its safe explicit-graph exponent pair
is $(C,\sigma)=(1,1)$, so
\[
 C+\sigma=2>2\cos(\pi/7)=1.8019377358\ldots.
\]
No cited result supplies a q-Horn implementation with the required smaller
joint exponent sum.
\end{lemma}
\begin{proof}
The explicit graph has $O(M)$ size; storing it and its traversal arrays uses
$O(M)$ work space.  This proves a safe upper bound, not a space lower bound.
The statement ``linear time'' does not imply sublinear writable space.
Implicit graph regeneration or low-space reachability can change both
resources, so it must be analyzed as a new same-realization algorithm rather
than combined with the published time bound.  Therefore only the displayed
conservative pair presently transfers.
\end{proof}

This is the selected highest-value negative result: adding a larger tractable
island is mathematically correct, but it does not yet satisfy the common
resource envelope needed by the final SAT algorithm.

\section{Concrete model and bounded test}
Code is
\path{experiments/proof_spines/qhorn_value_portfolio_test.py}; raw output is
\path{results/raw/proof_spines/2026-10-07_182659_qhorn_value_portfolio.json}.
The command
\begin{verbatim}
timeout 30s python experiments/proof_spines/qhorn_value_portfolio_test.py \
  results/raw/proof_spines/2026-10-07_182659_qhorn_value_portfolio.json
\end{verbatim}
uses seed 20261007 and a 25-second cooperative alarm.  It completed in
9.64 seconds with 30,048 checks.

The executable model enumerates the exact integer weights
$2\beta(x)\in\{0,1,2\}$ and independently checks truth tables.  The seeded
search examined 29,972 candidate formulas before finding
\[
\begin{split}
F={}&(\neg x_2\vee\neg x_3\vee\neg x_5)
\wedge(\neg x_2\vee x_3\vee\neg x_4)
\wedge(x_1\vee\neg x_3\vee\neg x_4)\\
&\wedge(x_1\vee x_3\vee x_4)
\wedge(x_1\vee x_3\vee x_5)
\wedge(x_1\vee x_4\vee x_5).
\end{split}
\]
The certificate is
$(2\beta(x_1),\ldots,2\beta(x_5))=(0,2,1,1,1)$.
Independent enumeration verifies that $F$ is satisfiable, not Horn,
dual-Horn, 2-CNF, or renamable Horn, not satisfied by either uniform phase,
and has no strong Horn backdoor of size at most two.  Thus the q-Horn member
has strict finite coverage beyond the currently implemented subportfolio.

The first precheck produced a four-variable q-Horn witness outside the four
named syntactic classes, but a size-two strong Horn backdoor still covered it.
That candidate was rejected rather than reported as a strict portfolio
witness; the test was strengthened to require five variables and absence of
every strong Horn backdoor of size at most two before the successful run.

Controls include a Horn chain, dual-Horn chain, 2-CNF cycle, contradictory
units, global complementation, three independent copies of the strict
witness, and tautological padding.  The code also checks the two resource
cells $(1,1)$ and hypothetical $(1,0)$ against the exact Williams constant.
The latter is explicitly hypothetical; the finite program does not prove it.

\section{Hostile, wrapper, and barrier audit}
\begin{description}
\item[] \textbf{Easy families:} Horn, dual-Horn, 2-CNF and renamable Horn
are absorbed by q-Horn; contradictory units confirm that the member must
return FALSE as well as TRUE.
\item[] \textbf{Renaming and padding:} signed complementation preserves
q-Horn after transforming the certificate.  Normalized tautologies do not
change the formula.
\item[] \textbf{Independent products:} variable-disjoint products remain
q-Horn by combining certificates; this gives a direct VALUE computation, not
a Cartesian backdoor cover.
\item[] \textbf{Current-scale/source carry:} the member never calls the
hypothetical $N^K$ SAT solver.  Its open issue is its own live workspace, not
source provenance.
\item[] \textbf{DEAD/COPY/DELAY/INTERP/RECOMP:} normalized dead clauses and
copies do not create coverage.  A wrapper may hide the q-Horn writing, in
which case recognition and access to the actual representation must be
charged anew.
\item[] \textbf{Oracle replacement:} a q-Horn or SAT oracle can reproduce
the returned bit.  The transfer relativizes and is not the nonrelativizing
hinge.
\item[] \textbf{Algebraic replacement:} algebraic access does not invalidate
the base-class upper bound and does not supply the missing all-CNF coverage.
No nonalgebrizing claim is made.
\item[] \textbf{Finite/asymptotic:} strict coverage of the displayed witness
is computational.  Linear-time recognition and decision come from the cited
theorems, not extrapolation from the test.
\end{description}

The post-quantum techniques in the standing ledger do not transfer here:
Schur squares, relation ranks, collision entropy, and random-oracle binding
do not yield a same-realization q-Horn work-space improvement.  Treating them
as relevant merely because q-Horn has a quadratic cover would be a
\textsc{Failed Transfer}.

\section{Cross-domain transfer classification}
The quadratic-cover construction maps q-Horn recognition to an implication
graph and strongly connected components.  This is an exact graph-algorithm
\textsc{Theorem Transfer}: the graph has linear size and the canonical
certificate is read from its SCC condensation.  What does not transfer is a
sublinear-space bound with the same linear time.  The missing bridge is an
explicit random-access implementation whose simultaneous time and work-space
exponents sum below $2\cos(\pi/7)$.

Control-theoretic Lyapunov functions and coding-theoretic low-rank
certificates are only analogies for resource accounting here; neither forces
such an implementation.  They are classified \textsc{Analogy/Candidate},
not theorem transfers.

\section{What changed relative to root main}
Root \texttt{main.tex} remains the unrelated LSPV synthesis and was not
edited.  Root \texttt{main5.tex} now states QHORN-VALUE1,
QHORN-SPACE-GATE1, the exact primary definition, the strict finite witness,
and the distinction between domain coverage and endpoint compatibility.
Root \texttt{main4.tex} records a concise no-status-change review: no
achievable pair or SAT-specific ARC closure is added.

\section{Adopted, rejected, and stale claims}
\textbf{Adopted:} q-Horn class containment; linear-time recognition and
decision; the Drive warning that q-Horn backdoor radius is a representation
statement; zero-child VALUE semantics; explicit acquisition and workspace
charging.

\textbf{Rejected/retired:} adding Horn, dual-Horn, 2-CNF, and renamable Horn
as four independent tractable coordinates; inferring a favorable joint
time--space pair from linear time alone; treating the exact-2P2N q-Horn
radius as a general SAT-machine lower bound.

\textbf{Stale or overstated claims found:} the current portfolio description
understated known q-Horn coverage and omitted the corresponding space gate.
No existing \textsc{Open} arrow was silently promoted.  QHORN-VALUE1 is exact
only as a partial base-class member.

\section{Open arrow and peer handoff}
The remaining target is
\[
\boxed{
\text{q-Horn recognition/decision with }C+\sigma<2\cos(\pi/7)
\quad\text{or a low-space selector that safely skips it,}
}
\]
embedded in an all-CNF favorable portfolio whose every success and failure
shares one charged envelope.  PROOF should audit whether the quadratic cover
can be queried implicitly with a rigorous time--space tradeoff; RED should
reject any proof that combines the published linear time with a different
low-space realization.  Even success on this target would still leave
universal favorable coverage open.  P versus NP remains open.

\bibliographystyle{plain}
\bibliography{spines}
\end{document}
% END_OF_VERIFIED_REPORT::RESEARCH::RSR-0004
